% Research contribution. This file is a section module, not a standalone document.
% Depends only on amsmath, amssymb, amsthm, hyperref and standard theorem environments.
\section{Published methods and exact improvements to the algebraic presentation}
\label{sec:lit-optimization}

Throughout this section a witness variable ranges over $\mathbb Q$ unless
another field is specified. An equivalence of polynomial presentations means
an equivalence of their projected rational solution sets in every displayed
input. Such an equivalence is independent of whether the arithmetic theorem
that identifies the original projected set with $\mathbb Z$ has been verified.
The results below should therefore be applied to the repository only after
matching their hypotheses to its actual residual system.

\subsection{An explicit three-variable semilocal definition}

Write
\[
 R=\mathbb Z_{(2)}\cap\mathbb Z_{(3)}
   =\{t\in\mathbb Q:v_2(t)\ge0,\ v_3(t)\ge0\},
\]
with $v_p(0)=+\infty$. The following specialization of classical
quaternion methods is particularly economical. Daans proves three-quantifier
existential definability for every finite intersection of valuation rings
in a global field; he traces the corresponding rational cases back to
Robinson \cite[Proposition 4.2 and Remark 4.3]{Daans2024}.
Here all constants and all exceptional cases are made explicit.

\begin{proposition}\label{prop:explicit-semilocal}
For every $t\in\mathbb Q$,
\[
 t\in R\quad\Longleftrightarrow\quad
 \exists u,v,w\in\mathbb Q:\quad
 u^2+6v^2-2w^2-3t^2+2=0.
\]
\end{proposition}
\begin{proof}
Consider the quaternion algebra $B=(2,3)_{\mathbb Q}$ with generators
$i^2=2$, $j^2=3$, $ij=-ji$. Its ramification set consists exactly of
$2$ and $3$. Indeed its Hilbert symbol at $3$ equals the Legendre symbol
$(2/3)=-1$, and at $2$ the formula
$(2,u)_2=(-1)^{(u^2-1)/8}$ for odd $u$ gives $(2,3)_2=-1$.
It splits at every other finite place because both parameters are units
at those odd primes, and it splits at
infinity because both parameters are positive. In particular $B$ is a
division algebra over $\mathbb Q$.

We use the standard quadratic embedding criterion for quaternion algebras:
a quadratic field $L/\mathbb Q$ embeds in $B$ if and only if every ramified
place of $B$ remains a field in $L$. This is the quadratic case of the
Albert--Brauer--Hasse--Noether theorem and the central simple algebra
embedding criterion; equivalently, extending scalars to $L$ kills the
local invariant $1/2$ precisely when the local extension has degree two.
See \cite[Chapter 14]{Voight2021}.

Suppose first that $t\ne0$, and put
\[
 \delta=1-\frac{2}{3t^2}.
\]
This number is nonzero: $t^2=2/3$ has no rational solution.
At $2$, put $m=v_2(t)$. If $m>0$, then $v_2(\delta)=1-2m$ is odd,
so $\delta$ is not a square in $\mathbb Q_2$. If $m=0$, then
$t^2\equiv1\pmod8$ and $\delta\equiv3\pmod8$, again a nonsquare.
If $m<0$, then $\delta\in1+8\mathbb Z_2$, hence it is a square in
$\mathbb Q_2$. At $3$, if $m=v_3(t)\ge0$ then
$v_3(\delta)=-1-2m$ is odd; if $m<0$, then
$\delta\in1+3\mathbb Z_3$ and Hensel's lemma makes it a square.
Consequently
\[
 \delta\notin\mathbb Q_2^{\times2},\quad
 \delta\notin\mathbb Q_3^{\times2}
 \quad\Longleftrightarrow\quad t\in R.
\]
The embedding criterion shows that this is equivalent to the existence
of a pure quaternion $xi+yj+zij$ with square $\delta$, that is,
\[
 2x^2+3y^2-6z^2=\delta.
\]
This equivalence also excludes the case in which $\delta$ is a rational
square: a division quaternion algebra cannot contain a noncentral element
whose square is a nonzero rational square, since
$(\alpha-s)(\alpha+s)=0$ would force $\alpha=\pm s$.
A nonzero pure quaternion has reduced trace zero, so it cannot be such
a nonzero rational scalar.
Multiplying the displayed equation by $3t^2$ and setting
$u=3ty$, $v=tx$, $w=3tz$ gives the equation in the statement.
These substitutions are invertible for $t\ne0$.
Finally $t=0$ belongs to $R$ and is represented by $(u,v,w)=(0,0,1)$.
\end{proof}

\begin{lemma}\label{lem:two-semilocal-inputs}
Let $F(A,B)=A^2+AB-B^2$. For every $a,b\in\mathbb Q$ and $p=2,3$,
\[
 v_p(F(a,b))=2\min\{v_p(a),v_p(b)\}.
\]
In particular $(a,b)\in R^2$ if and only if $F(a,b)\in R$.
\end{lemma}
\begin{proof}
The assertion is immediate for $a=b=0$. Otherwise scale $a,b$ by a
common power of $p$ until both are integral and at least one is a unit.
The reduction of $A^2+AB-B^2$ has no nontrivial zero over either
$\mathbb F_2$ or $\mathbb F_3$: over $\mathbb F_2$ one checks the
three nonzero pairs, and over $\mathbb F_3$ the discriminant $5$
is a nonsquare. Thus the scaled value is a unit, proving the equality.
\end{proof}

\begin{corollary}[Five-variable replacement for the anchor]
\label{cor:five-anchor}
Fix rational inputs $f_0,f_1,C$ with $C\ne0$. The condition
$f_0/C,f_1/C\in R$ is equivalent to the existence of five rational
variables $r,D,u,v,w$ satisfying the three quadratic equations
\begin{align*}
 r&=f_0^2+f_0f_1-f_1^2,\\
 D&=C^2,\\
 u^2+6v^2-2w^2-3r^2+2D^2&=0.
\end{align*}
\end{corollary}
\begin{proof}
By Lemma~\ref{lem:two-semilocal-inputs}, the desired condition is
$r/C^2\in R$, with $r=F(f_0,f_1)$. Apply
Proposition~\ref{prop:explicit-semilocal} to $t=r/C^2$ and multiply
its equation by $C^4$. Multiplication of each of its three witnesses
by $C^2$ is a bijection of $\mathbb Q$, giving the stated equation.
\end{proof}

The released anchor uses two scaling variables and two copies of a
seven-witness trace-sum condition, hence sixteen witnesses in this module.
Corollary~\ref{cor:five-anchor} replaces that module by five witnesses.
It therefore saves eleven variables before any other change and retains
quadratic residuals, so their sum of squares remains quartic. This count
concerns the specified module; an implementation should report the exact
final global count after removing any unused lifting definitions.
The three displayed residuals have integer coefficients of absolute value
at most six. This is an exact replacement, with no change to the arithmetic
meaning or strength of the anchor.

The same method handles arbitrary finite tuples. For a finite nonempty
set $S$ of rational primes and a positive integer $m$, choose a monic
polynomial $g\in\mathbb Z[T]$ of degree $m$ whose reductions are
irreducible modulo every $p\in S$; finite-field irreducible polynomials
and the Chinese remainder theorem supply one. If $\alpha$ is a root, let
\[
 N_m(X_1,\ldots,X_m)=
 N_{\mathbb Q(\alpha)/\mathbb Q}
 (X_1+X_2\alpha+\cdots+X_m\alpha^{m-1}).
\]
For $p\in S$ the extension is unramified and inert at $p$, and its
power basis is an integral basis locally. Therefore
\[
 v_p(N_m(\mathbf x))=m\min_i v_p(x_i).
\]
This proves, with no new existential variables, that all coordinates of
$\mathbf x$ belong to $\bigcap_{p\in S}\mathbb Z_{(p)}$ if and only if
$N_m(\mathbf x)$ does. It is the norm-form version of
\cite[Proposition 4.8]{Daans2024}. In the present two-prime, two-input
case the polynomial $T^2-T-1$ gives precisely the form $F$ above.

\subsection{Precision thresholds and a smaller scalar order}

The following elementary argument isolates exactly what the growth
comparison needs. It does not establish the arithmetic estimates from
which it starts.

\begin{proposition}\label{prop:precision-thresholds}
Let $a,B,n,\Delta\in\mathbb Z$, with $B\ge a\ge1$, let
$\beta\in\mathbb Q$, and let $H\ge1$ be real.
Write $h(\beta)=\log\max\{|b|,c\}$ when $\beta=b/c$ is in lowest
terms with $c>0$. If $\beta\ne n$, let $N$ be the absolute value
of the numerator of $\beta-n$ in lowest terms.
Suppose an arithmetic certificate supplies the following estimates
for some $\kappa>48$,
$\lambda>12$, and $\mu>24$:
\begin{enumerate}
\item if $\Delta\ne0$, then
$B^2-6<7H|\Delta|^{24/\kappa}$ and $|\Delta|<5B^4$;
\item if $\Delta=0$, then $\log B>a-3$;
\item if $\beta\ne n$, then
$B^2<7HN^{24/\lambda}+6$ and $N\le2e^{h(\beta)}B$;
\item $h(\beta)\le H a^{24/\mu}$.
\end{enumerate}
Put
\[
 C_\kappa=7\,5^{24/\kappa}+6,\qquad
 C_\lambda=7\,2^{24/\lambda}+6,\qquad
 A_\lambda=\frac{\lambda}{2\lambda-24},\qquad
 B_\lambda=\frac{12}{\lambda-12}.
\]
If
\[
 a>\max\left\{6,
 (C_\kappa H)^{\kappa/(2\kappa-96)},
 4A_\lambda\log(C_\lambda H),
 (4B_\lambda H)^{\mu/(\mu-24)}\right\},
\]
then $\Delta=0$ and $\beta=n$.
\end{proposition}
\begin{proof}
For $\Delta\ne0$, absorb the constant six using $B,H\ge1$ to obtain
$B^{2-96/\kappa}<C_\kappa H$, contradicting $B\ge a$.
Thus $\Delta=0$. If $\beta\ne n$, put $r=24/\lambda<2$.
The third hypothesis and $B,H\ge1$ imply
\[
 B^{2-r}<C_\lambda H e^{r h(\beta)}.
\]
Consequently
\[
 \log B<A_\lambda\log(C_\lambda H)+B_\lambda h(\beta)<a/2<a-3,
\]
contrary to the second hypothesis.
\end{proof}

The thresholds $48,12,24$ are thresholds of this particular comparison
of exponents. This proof does not show that the corresponding boundary
values are impossible by a different argument.

\begin{corollary}\label{cor:scalar32}
The choices
\[
 (\kappa,\lambda,\mu)=(52,13,32),\qquad a>(21H)^{13/2}
\]
satisfy the conclusion of Proposition~\ref{prop:precision-thresholds}.
If $\log_2H<4783227$, the existing type of anchor with twenty-six
successive squarings, giving $a\ge3^{2^{25}-1}$, is sufficient.
\end{corollary}
\begin{proof}
Here $C_\kappa<21$, $C_\lambda<32$,
$A_\lambda=13/2$, $B_\lambda=12$, and $24/\mu=3/4$.
The assumed lower bound for $a$ implies $a>(48H)^4$ because
$21^{13/2}>48^4$ and $H\ge1$; hence $h(\beta)<a/48$.
It also implies $a>26\log(32H)$: use $\log(32H)<32H$ and
$21^{13/2}>832$. Thus the last contradiction reads
$\log B<(13/2)\log(32H)+12h(\beta)<a/2<a-3$.
For the numerical anchor,
\[
 \log_2\big((21H)^{13/2}\big)
 <\tfrac{13}{2}(5+4783227)=31091008,
\]
whereas
\[
 \log_2\big(3^{2^{25}-1}\big)
 >\tfrac32(2^{25}-1)=50331646.5.
\]
\end{proof}

For the source system's scalar batch the order $48$ is used solely to
obtain $M_5(\beta)^{48}\mid s$, hence
$h(\beta)\le H\sqrt{|s|}$. Replacing that order by $32$ gives instead
$h(\beta)\le H|s|^{3/4}$, which Corollary~\ref{cor:scalar32} shows
is sufficient. Completeness still allows all the required divisibilities:
the index $s$ is chosen in an arbitrarily divisible nonzero subgroup.
The batch's recursive power construction uses eight multiplications
for the required powers $r^2,r^{96}$, and six for $r^2,r^{64}$.
Thus this change removes two source power definitions. Its exact effect
on the compiled polynomial must be measured after the full pipeline.

For comparison, orders $49$ and $52$ require the same thirteen terms
in the logarithm truncation used by the source, but the source's shared
power recursion is less economical for $(49,13)$: it uses twelve
multiplications for $r^2,r^{26},r^{98}$ and nine for
$r^2,r^{26},r^{104}$. A smaller numerical precision does not automatically
mean fewer variables. Joint addition-chain optimization is the correct
comparison.

\subsection{Combining square conditions without assuming integrality}

Matiyasevich and Robinson's relation-combining theorem is an essential
source of small integer Diophantine representations
\cite{MatiyasevichRobinson1975,Jones1982}. Its integer-domain hypotheses
must not be silently replaced by rational-domain hypotheses. The
appropriate rational version is Zhang--Sun
\cite[Theorem 1.2]{ZhangSun2022}. Their separation parameter is
$(k+\sum A_i^2)(1+\sum A_i^{-2})$, and their theorem assumes that every
$A_i$ is nonzero. The following slightly simpler separation bound also
works and has a direct proof.

\begin{lemma}\label{lem:additive-square-separation}
Let $k\ge1$. For nonzero $A_1,\ldots,A_k\in\mathbb Q$, put
\[
 W=\sum_{i=1}^k(A_i^2+A_i^{-2}),\qquad
 J(T)=\prod_{\boldsymbol\epsilon\in\{\pm1\}^k}
 \left(T-\sum_{i=1}^k\epsilon_i W^{i-1}\sqrt{A_i}\right).
\]
Then $J\in\mathbb Q[T]$, and $J$ has a rational root if and only if
all $A_i$ are rational squares. Every complex root of $J$ is nonzero.
\end{lemma}
\begin{proof}
Independently changing the formal sign of each square root fixes the
product, so its coefficients are rational polynomials in the $A_i$ and
$W$. Let
$m=\min_i\sqrt{|A_i|}$, $M=\max_i\sqrt{|A_i|}$, and $\rho=M/m\ge1$.
The two summands corresponding to maximal and minimal absolute value
show that $W\ge2\rho^2\ge1+\rho$; the same inequality is immediate
when both extrema occur at one index. Therefore for every $j$,
\[
 \sum_{i<j}M W^{i-1}<\frac{M W^{j-1}}{W-1}
 \le mW^{j-1}.
\]
No nonempty signed subsum of the radicals can vanish, since its last
term has larger absolute value than the sum of the preceding terms.
In particular every root is nonzero. If a displayed root is rational,
every automorphism of the multiquadratic splitting field fixes it.
Subtracting its conjugate would give a nonempty signed subsum if any
radical changed sign. Thus every radical is fixed by the full Galois
group and is rational. The converse is immediate.
\end{proof}

There is also an explicit one-variable polynomial that enforces the
nonzero condition itself. Put
\[
 S=\prod_{i=1}^k A_i,\qquad D=S^2,\qquad
 N=D\sum_i A_i^2+\sum_i\prod_{j\ne i}A_j^2
\]
and define
\[
 K_k(\mathbf A,U)=\prod_{\boldsymbol\epsilon\in\{\pm1\}^k}
 \left(1-U\sum_{i=1}^k\epsilon_i
 S^{2k-2i+1}N^{i-1}\sqrt{A_i}\right).
\]
Independent sign changes show that this is an integer-coefficient
polynomial in $\mathbf A,U$. If $S=0$ it is identically $1$.
If $S\ne0$, the signed sum in each factor is
$S^{2k-1}\sum_i\epsilon_iW^{i-1}\sqrt{A_i}$, so
Lemma~\ref{lem:additive-square-separation} proves
\[
 \exists U\in\mathbb Q:K_k(\mathbf A,U)=0
 \quad\Longleftrightarrow\quad
 A_i\in\mathbb Q^{\times2}\text{ for all }i.
\]
This polynomial has degree $2^k$ in $U$; total degree grows
exponentially with $k$. To allow zero square inputs, take the finite
product, over $I\subseteq\{1,\ldots,k\}$, of
\[
 \sum_{i\notin I}A_i^2+K_{|I|}((A_i)_{i\in I},U)^2,
\]
using zero in place of $K_0$ for the empty subset. One shared rational
$U$ suffices for all branches. This treats zero inputs exactly, but its
expanded degree is correspondingly large.

Private square roots can be replaced in this way only when no other
condition needs their particular values or signs. If roots occur elsewhere,
the complete finite algebraic system must be compressed, as below.
For example, replacing $y^2=A,\ y=1$ by the bare condition that $A$
is a square would be invalid.

\subsection{One variable for finite choices and nonzero conditions}

The useful general principle is that the number of finite algebraic
choices need not equal the number of rational witnesses. Sun's recent
preprint \cite[Sections 10.3--10.6]{Sun2026Rational86} uses this principle
for a bounded-variable undecidability construction. The elementary
algebra below is independently verifiable and does not depend on that
preprint's height or prime-pattern estimates. Its undecidability result
is not itself a definition of $\mathbb Z$ inside $\mathbb Q$.

\begin{lemma}\label{lem:finite-separation}
Let $A$ be a reduced finite-dimensional $\mathbb Q$-algebra of dimension
$d\ge1$, generated by $N\ge1$ elements $a_1,\ldots,a_N$. Some integer
\[
 1\le b\le1+(N-1)\binom d2
\]
makes $z_b=a_1+ba_2+\cdots+b^{N-1}a_N$ have distinct values under all
$d$ embeddings $A\to\overline{\mathbb Q}$. For such a $b$, the
characteristic polynomial of multiplication by $z_b$ has a rational root
if and only if $A$ has a $\mathbb Q$-algebra homomorphism to $\mathbb Q$.
\end{lemma}
\begin{proof}
Two different embeddings differ on a generator, so equality of their
values on $z_b$ is a nonzero polynomial equation in $b$ of degree at most
$N-1$. At most $(N-1)\binom d2$ values are bad. For a good value the
characteristic polynomial has precisely the distinct embedding values
as its roots. If one value is rational, the absolute Galois group fixes
that value and hence its unique corresponding embedding; the values
of all generators are therefore rational. Conversely a rational
homomorphism supplies a rational root.
\end{proof}

\begin{lemma}\label{lem:reciprocal-guard}
Suppose
\[
 P_{\mathbf x}(T)=T^d+\sum_{j=1}^d
 \frac{N_j(\mathbf x)}{B(\mathbf x)}T^{d-j}
\]
is monic on $B\ne0$, and let $H(\mathbf x)$ be the product of the
numerators of all further required inequations. Then
\[
 \Psi(\mathbf x,U)=1+\sum_{j=1}^d
 N_j(\mathbf x)B(\mathbf x)^{2j-1}H(\mathbf x)^jU^j
\]
has a rational root in $U$ exactly when $BH\ne0$ and $P_{\mathbf x}$
has a nonzero rational root.
\end{lemma}
\begin{proof}
If $BH=0$, all nonconstant terms vanish and $\Psi=1$.
If $BH\ne0$, set $L=B^2H$. For $U\ne0$ the identity
\[
 \Psi(\mathbf x,U)=(LU)^dP_{\mathbf x}(1/(LU))
\]
proves the assertion, and $\Psi(\mathbf x,0)=1$ excludes $U=0$.
\end{proof}

Zero roots are covered by taking the two finite alternatives $P(T)$
and $P(T-1)$: a rational root of one of them is nonzero. Thus nonzero
conditions need no additional inverse variable once one finite-choice
variable is already present. Equations depending only on $\mathbf x$
can be imposed by adding their squares to $\Psi^2$, and finite alternatives
can be combined by multiplying the resulting polynomials, using the
same $U$ in every factor.

For completeness, this also handles nonreduced finite algebras and
additional equations involving their generators. Quotient by the ideal
of those equations, and discard factors violating an inequation by
localizing at its value; these are finite-dimensional algebra operations.
Pass to the reduced quotient. In characteristic zero its nilradical is
the kernel of the trace pairing $(a,b)\mapsto\operatorname{Tr}(ab)$:
a product decomposition into Artinian local factors reduces the assertion
to the nondegeneracy of the trace pairing of a separable field extension.
Gaussian elimination then computes this quotient. When coefficients depend
on free parameters, branch on the finitely many pivots being zero or
nonzero and record all denominators. On every branch apply
Lemma~\ref{lem:finite-separation}, enforce the nonzero discriminant of
the selected characteristic polynomial, and apply
Lemma~\ref{lem:reciprocal-guard}. This is an effective finite construction
when the original algebra has a fixed finite basis, for example for a
tower of monic equations. It preserves exceptional specializations.

Dropping the separation condition is unsound: $x^2=y^2=2$ has no
rational pair, but the characteristic polynomial for $x+y$ is
$T^2(T^2-8)$ and has the rational root zero. Its repeated roots signal
precisely the collision. Using $x+2y$ instead gives
$T^4-20T^2+36$, with no rational root.

\subsection{How MRDP transfers, and what it does not compress}

Let $P(t,\mathbf y)$ be an integer-coefficient definition of $\mathbb Z$
in $\mathbb Q$ with $n$ witnesses and degree $d$. For an arbitrary
integer polynomial $f(X_1,\ldots,X_m)$, define
\[
 Q_f(\mathbf X,\mathbf Y_1,\ldots,\mathbf Y_m)=
 f(\mathbf X)^2+\sum_{i=1}^mP(X_i,\mathbf Y_i)^2.
\]
Then $f$ has an integer zero if and only if $Q_f$ has a rational zero.
This uses $m(n+1)$ rational variables and has degree
$2\max\{\deg f,d\}$. Both directions follow immediately from the
nonnegativity of rational squares. The transformation is effective.

MRDP yields more than undecidability after this transfer. If
$S\subseteq\mathbb Q$ is computably enumerable, the integer pairs
$(a,b)$ with $b>0$, $\gcd(a,b)=1$ and $a/b\in S$ form a computably
enumerable subset of $\mathbb Z^2$. MRDP supplies an integer polynomial
$F(a,b,\mathbf z)$ defining those pairs using integer witnesses.
Apply $P$ to every integer variable in $F$ and add the equation $bq-a=0$.
The resulting rational existential formula defines $S$.
Consequently, conditional on $P$, the computably enumerable subsets of
$\mathbb Q$ are exactly the Diophantine subsets of $\mathbb Q$.
The reverse implication is effective enumeration of rational witness tuples.
The same argument uniformly in an index gives a universal rational
Diophantine polynomial.

This conclusion does not permit us to use the small integer universal
polynomials directly with rational witnesses. Nor can a merely
integer-valued Diophantine model of arithmetic be identified with the
ordinary embedded integers without a separate value relation. Universal
integer formulas, a Diophantine model, an undecidability reduction, and
an existential definition of the actual subset $\mathbb Z\subset\mathbb Q$
are distinct statements.

More generally finite unions allow complete reuse of witness coordinates:
if $S_j$ has a presentation $\exists\mathbf y_j\,f_j(t,\mathbf y_j)=0$,
then their union has the presentation
\[
 \exists(y_1,\ldots,y_N):
 \prod_j f_j(t,y_1,\ldots,y_{n_j})=0,
 \qquad N=\max_j n_j.
\]
Conjunctions do not in general allow that reuse. Daans--Dittmann--Fehm
prove an improved existential-rank intersection theorem over fields finitely
generated over a perfect subfield, and Daans states the usable form
\cite[Theorem 2.5]{Daans2024}: ranks $m_1,m_2\ge1$ give intersection
rank at most $m_1+m_2-1$. For a concrete polynomial, an explicit
transformation such as the finite-algebra construction above is preferable
to claiming an unimplemented numerical saving from that abstract theorem.

Finally, rational nonnegativity cannot be replaced by three ordinary
squares: $7$ is not a sum of three rational squares. Clearing a common
denominator would give $7q^2=a^2+b^2+c^2$; after removing common powers
of two, the standard mod-eight obstruction contradicts this equation.
Lagrange's four-square theorem does give every nonnegative rational as a
sum of four rational squares, by representing the product of numerator
and denominator and dividing by the denominator squared. General nonlinear
three-witness definitions require a separate argument.

To compare coefficient sizes, first clear numerical denominators and
normalize the polynomial to primitive integer content. Rational invertible
changes of witness coordinates and nonzero constant rescaling preserve the
projection, but may change coefficient height. Replacing large constants
by addition or multiplication chains decreases coefficient size at the cost
of new variables; under the requested priority of minimizing variables,
such a trade is secondary to exact variable elimination.
